\documentclass[UTF8,11pt]{ctexart}

\usepackage[a4paper,margin=2.3cm]{geometry}
\usepackage{amsmath,amssymb,bm}
\usepackage{booktabs,array}
\usepackage{xcolor}
\usepackage{hyperref}
\hypersetup{colorlinks=true,linkcolor=blue,urlcolor=blue}

\newcommand{\Vp}{V_\mathrm{p}}
\newcommand{\Sw}{S_\mathrm{w}}
\newcommand{\Sp}{S_\mathrm{p}}
\newcommand{\Aflux}{A_\mathrm{flux}}
\newcommand{\Lax}{L_\mathrm{ax}}
\newcommand{\Pfus}{P_\mathrm{fus}}
\newcommand{\Pbrem}{P_\mathrm{brem}}
\newcommand{\Pcycl}{P_\mathrm{cycl}}
\newcommand{\Pheat}{P_\mathrm{heat}}
\newcommand{\Pwall}{P_\mathrm{wall}}
\newcommand{\Eth}{E_\mathrm{th}}
\newcommand{\iotabar}{\bar{\iota}}

\title{PolyFusion 仿星器几何：怎么画、怎么进功率账}
\author{PolyFusion release —— 仿星器几何系列报告（配 docs/39）}
\date{\today}

\begin{document}
\maketitle

\begin{abstract}
本报告说明 PolyFusion 0-D 程序里仿星器的几何\textbf{当前到底怎么实现}，以及这套几何\textbf{以哪些量、按什么幂次进入功率账}。核心结论：几何有\textbf{两个互相独立又互相一致}的角色——(1) 前端\textbf{截面视图}（cartoon，给人看形状），(2) 后端\textbf{功率账几何}（一组标量 $a,\Vp,\Sw,\Lax,\iotabar$，给功率公式做积分权重）。两者通过\textbf{同一个截面面积锚} $\Aflux=\pi a^2$ 绑定，已数值核对一致到 $\le 1\%$。
\end{abstract}

\section{两个几何角色}

\paragraph{(A) 功率账几何（后端，真正进功率）。}
功率公式只需要一小组\textbf{标量}：等离子体体积 $\Vp$、第一壁面积 $\Sw$、面积等效小半径 $a$、大半径 $R_0$、磁轴长 $\Lax$、旋转变换 $\iotabar$。它们不依赖"画成什么形状"。

\paragraph{(B) 截面视图几何（前端，只给人看）。}
UI 在 $\varphi=0,\,T/4,\,T/2$ 三个环向切面画磁面截面 + 壁层。这是\textbf{装饰性 cartoon}，不进任何功率数。

二者的唯一物理绑定：\textbf{视图截面的面积 $=$ 功率账用的 $\Aflux=\pi a^2$}（见 \S\ref{sec:consistency}）。

\section{功率账几何怎么取（后端）}\label{sec:backend}

固定关系：
\begin{equation}
a = \frac{R_0}{A},\qquad \Aflux = \pi a^2 ,
\end{equation}
其中 $A$ 是环径比（aspect ratio）输入。$\Aflux$ 取\textbf{一阶近轴的通量守恒截面积}：一阶近轴椭圆拉长比 $\kappa$ 与 $1/\kappa$ 相乘抵消，面积恒为 $\pi a^2$（面积守恒），所以\textbf{拉长比不改体积}。

体积/壁面按位形类别两条路：
\begin{itemize}
\item \textbf{概念堆}（HELIAS、NAE-QA，本就是近轴设计）：用近轴几何
  \begin{equation}
    \Vp = \Aflux\,\Lax = \pi a^2\,\Lax,\qquad
    \Sw = \bar{C}_\mathrm{wall}\,\Lax,
  \end{equation}
  $\Lax$ 为近轴磁轴弧长，$\bar{C}_\mathrm{wall}$ 为弧长加权的壁截面周长（边界 $+g$）。$\iotabar$ 由周期 $\sigma$ 方程自洽解出（含轴扭率贡献）。
\item \textbf{真机}（W7-X、LHD、HSX、CFQS）：单谐波近轴\textbf{无法}表示，改用\textbf{实测覆写}
  \begin{equation}
    \Vp = \Vp^\mathrm{override},\quad
    \Sw = \Sw^\mathrm{override},\quad
    \iotabar = \iotabar^\mathrm{measured}.
  \end{equation}
  此时近轴解被丢弃；其 $\iotabar_\mathrm{geom},\kappa_\mathrm{eff},e_\mathrm{max}$ 仅作\textbf{诊断估计}（UI 标"\,$\sim$近轴估"），并把上报的 $V_\mathrm{p,geom},S_\mathrm{w,geom}$ 置为实际所用值以免误导。
\end{itemize}

\section{几何怎么进功率账（哪个量、什么幂次）}\label{sec:power}

总功率平衡（与其它位形共用结构）：
\begin{equation}
\Pheat = \Pcycl + \Pbrem + P_\mathrm{line} + \frac{\Eth}{\tau_E} - f_\mathrm{ion}\Pfus .
\end{equation}
几何标量的进入点如下表。

\begin{center}
\renewcommand{\arraystretch}{1.35}
\begin{tabular}{@{}lll@{}}
\toprule
功率量 & 公式（几何相关因子）& 几何依赖 \\
\midrule
聚变功率 & $\Pfus \propto n_{10}n_{20}\,\Phi\,\Vp$ & $\propto \Vp$（体积积分权重）\\
中子功率 & $P_n=\Pfus(1-f_\mathrm{ion})$ & 经 $\Pfus$ 经 $\Vp$ \\
轫致辐射 & $\Pbrem \propto n_e^2\sqrt{T_e}\,\Vp$ & $\propto \Vp$ \\
回旋辐射 & $\Pcycl \propto B_0^{2.5}\,a^{-0.5}\,\Vp$ & $\propto \Vp\,a^{-1/2}$ \\
储能/输运 & $\Eth \propto (n_iT_i+n_eT_e)\,\Vp$,\ \ $P_\mathrm{tr}=\Eth/\tau_E$ & $\propto \Vp$ \\
第一壁负荷 & $\Pwall = (\Pfus+\Pheat)/\Sw$ & $\propto 1/\Sw$ \\
ISS04 约束 & $\tau_\mathrm{ISS04}\propto a^{2.28}R_0^{0.64}P_L^{-0.61}\bar{n}^{0.54}B_0^{0.84}\iotabar^{0.41}$ & $a,R_0,\iotabar$ \\
Sudo 密度限 & $n_\mathrm{Sudo}=0.25\sqrt{P_L B_0/(a^2R_0)}\cdot10^{20}$ & $a,R_0$ \\
\bottomrule
\end{tabular}
\end{center}

要点：
\begin{itemize}
\item \textbf{体积 $\Vp$ 是最强权重}：$\Pfus,\Pbrem,\Pcycl,\Eth$ 全部线性正比于 $\Vp$。真机用实测 $\Vp$，所以聚变/辐射功率锚在实测体积上。
\item \textbf{$\Sw$ 只进壁负荷}：$\Pwall=(\Pfus+\Pheat)/\Sw$。这就是真机要 $\Sw^\mathrm{override}$ 的原因——近轴 $S_\mathrm{w,geom}$ 因极端拉长而失真（如 HSX 近轴 $e_\mathrm{max}\!\sim\!384$）。
\item \textbf{$a,R_0,\iotabar$ 进约束闭合}：ISS04 约束时间与 Sudo 密度限。$\iotabar$ 仅在此以 $0.41$ 次幂出现；真机用实测 $\iotabar$。
\item \textbf{形状（bean / 旋转椭圆 / D 形）不进功率}：功率只看上述标量。截面形状纯属可视化。
\end{itemize}

\section{截面视图怎么画（前端）}\label{sec:view}

\paragraph{概念堆——二阶近轴。}
用 Garren–Boozer / Landreman–Sengupta–Plunk 的\textbf{二阶}（$r^2$）近轴展开（移植 pyQSC，对齐 $<2\times10^{-12}$）。截面
\[
P(\theta)=\text{axis}+X\hat n+Y\hat b+Z\hat t,\quad
X=rX_{1c}\cos\theta+r^2(X_{20}+X_{2c}\cos2\theta+X_{2s}\sin2\theta),\ \dots
\]
一阶体按真实 $a$ 画（尺寸对），二阶 wobble 限幅（$\le$ 一阶的 $50\%$）给出 bean/crescent，互不自交。

\paragraph{真机——真实公开平衡边界谐波。}
取 DESC 公开平衡（\texttt{desc/examples}：W7-X、HSX、ESTELL；经典 $l{=}2$ heliotron 作 LHD）的边界 Fourier 系数，截断 $|m|,|n|\le2$、归一化（减 $R_{00}$、除原生小半径），按 DESC 双 Fourier 乘积基求值并重缩放到预设 $R_0,a$：
\begin{equation}
R(\theta,\varphi)=R_0+a\!\sum_{m,n}\!R_{mn}A_m(\theta)B_n(\varphi),\quad
Z(\theta,\varphi)=a\!\sum_{m,n}\!Z_{mn}A_m(\theta)B_n(\varphi),
\end{equation}
$A_m=\cos m\theta\,(m\!\ge\!0)|\sin|m|\theta\,(m\!<\!0)$，$B_n=\cos(n\,n_\mathrm{fp}\varphi)|\sin(\dots)$。这复现真实的 W7-X bean$\to$triangle、LHD 旋转椭圆、HSX/CFQS 形态。

\paragraph{嵌套面与壁层。}
嵌套磁面与壁层都\textbf{从该切面真实质心}按比例缩放（$n\!\ne\!0$ 谐波会把每个切面挪离 $R_0$，若从 $R_0$ 缩放会把磁面甩到壁外）。保证 磁面 $\subset$ 边界 $\subset$ 壁。UI 提供\textbf{相位开关}（默认全显，亦可单看 $\varphi=0/T4/T2$）。

\section{一致性核查}\label{sec:consistency}

对全部 6 预设数值核对（前端截面面积 vs 后端 $\Aflux$；前端隐含体积/壁面 vs 后端 $\Vp/\Sw$）：

\begin{center}
\renewcommand{\arraystretch}{1.25}
\begin{tabular}{@{}lcccc@{}}
\toprule
预设 & $a$ & $\Aflux=\pi a^2$ & 前端截面面积 & 比值 \\
\midrule
HELIAS / NAE-QA & 1.80 & 10.18 & 10.1 & 0.99–1.00 \\
W7-X & 0.55 & 0.950 & 0.946 & 1.00 \\
LHD  & 0.65 & 1.327 & 1.327 & 1.00 \\
HSX  & 0.15 & 0.071 & 0.070 & 0.99 \\
CFQS & 0.25 & 0.196 & 0.196 & 1.00 \\
\bottomrule
\end{tabular}
\end{center}

\begin{itemize}
\item \textbf{核心一致}：前端 $a,R_0,n_\mathrm{fp}$ $=$ 后端；\textbf{前端截面面积 $=$ 后端 $\Aflux=\pi a^2$（$\le1\%$，全部预设）}。真机谐波按原生小半径归一化使面积恒为 $\pi a^2$。
\item \textbf{概念堆}：前后端同一近轴模型，体积/壁面一致到 $1$–$5\%$。
\item \textbf{真机}：功率账用实测 $\Vp/\Sw$；前端真实形状的几何体积/壁面在其 $\sim10$–$25\%$ 内（W7-X $32.7$ vs $30$；$\Sw$ $141$ vs $128$），实测值为准。
\item 真机近轴诊断（$\iotabar_\mathrm{geom},\kappa_\mathrm{eff},e_\mathrm{max}$）已在 UI 标"$\sim$近轴估"，不进功率数。
\end{itemize}

\section{诚实局限}
\begin{itemize}
\item 截面视图是 cartoon：真机用 $|m|,|n|\le2$ 截断 + 归一化重缩放，非完整平衡重建；LHD 用通用 $l{=}2$ heliotron 结构、CFQS 用 ESTELL 代理（形态族正确、非逐机精确）。
\item 真机 $\Vp/\Sw/\iotabar$ 覆写值为公开估计，需对照权威文献核实。
\item 嵌套面 $\rho$ 线性自相似、壁层比例外扩，均为显示近似。
\end{itemize}

\paragraph{来源。}DESC \texttt{desc/examples}（github.com/PlasmaControl/DESC）；ISS04：Yamada NF 45 (2005) 1684；Sudo NF 30 (1990)；近轴：Landreman–Sengupta–Plunk JPP 85 (2019)。代码细节见 \texttt{polyfusion/configs/stellarator.py}、\texttt{polyfusion/nearaxis.py}、\texttt{presets/stellarator.json}，过程见 \texttt{docs/39}。

\end{document}
